\section{Implementation and deployment}
\label{sec:impl}

The platform is implemented as a single server runtime in TypeScript on
Node.js, with a Vue 3 front end. One process hosts the four entry points
(HTTP, WebSocket, A2A, MCP), the agent runtime, and the industrial stack;
there is no separate governance service to deploy or bypass. State lives
in an embedded SQLite database (via the Node built-in driver) by default;
when production-grade history is required, a provisioned deployment adds a
TimescaleDB (compressed, chunked daily) and an MQTT broker, and the
acquisition gateway switches its storage backend without changes to the
governed path.

\textbf{Protocol surface.} Five industrial connectivity protocol drivers
are built in:
Modbus TCP, Modbus RTU, OPC UA \cite{iec62541}, MQTT, and HTTP, plus a
serial-port plugin driver and a MES REST integration with mirror and live
modes. Drivers implement a common read/write contract, so admission
limits, approval, and receipt grading apply uniformly; a driver cannot
weaken them. Write nodes carry per-driver configuration (register maps,
byte order, scaling, topic maps) declared once in the scene configuration.

\textbf{Agent harness surface.} Fourteen harness adapters execute agent
teams: two in-process engines (a deterministic mock for testing and a
Claude-compatible engine), an RPC-subprocess engine, an app-server engine,
three ACP-style engines, a served engine, and six headless CLI adapters.
The registry normalizes turn execution and tool invocation, so the
governed operation surface is identical across engines; an episode can mix
engines per channel. Teams are described declaratively (lead and worker
roles, model preferences) and instantiated from templates.

\textbf{Configuration and extension.} A single configuration schema feeds
defaults (a YAML file), runtime overrides, and environment variables;
settings and their descriptions are declared once and consumed by both the
settings UI and the CLI. A hot-reloadable plugin surface registers
additional drivers, processors, node templates, agent tools, HTTP routes,
browser panels, twin providers, and line-permission hooks; a failing
plugin is isolated from the runtime. The platform ships as an installable
npm package with an \texttt{aw} CLI (\texttt{aw start} initializes a home
directory, migrates prior versions, and serves the UI); production
posture terminates TLS on a reverse proxy in front of the single service.
The host tool surface exposed to agents comprises 105 built-in tools plus 21
plugin-registered tools (126 at the evaluated release), including the
industrial operations used in this paper; every tool's description and
scope semantics are generated from the same operation catalog that the
governance block enforces.

\textbf{Evidence pipeline.} All experiments in Section~\ref{sec:eval} are
driven by a benchmark pipeline in the same repository. A runner provisions
scenes through the platform's own REST surface, executes scenario scripts,
and freezes a report, raw event logs, and checksums into a versioned
results directory. Each result records the runner command, source commit,
and harness hash, which is what enables the frozen-baseline reproduction
checks of Section~\ref{sec:eval-quality}.

\textbf{Deployment posture and its boundaries.} The single-process design
places the agent runtimes (including outbound-capable harness adapters) and
the protocol write drivers in one trust domain. That is a deliberate
simplicity for evaluation, but it does not by itself satisfy IEC 62443-3-2
zoning. A plant deployment should place the platform inside a defined zone
with conduits to the control network, isolate harness adapters from
field-facing interfaces, and provide backup and restart procedures. Node
version, resource sizing, and high-availability specifications are
deployment-engineering tasks this paper does not claim; harness-side
isolation, air-gapped packaging, and a zoned reference architecture are
future work (Section~\ref{sec:conclusion}).
